\begin{definition}[Local tile dynamics bridge status]
\label{def:philosophy-boundary-tile-bridge-status}
The local-tile dynamics bridge
$\mathsf{LocalTileDynamics}^{\uparrow}$ is a
$\mathsf{VisionSeed}/\mathsf{BiologicalBridge}$ target and a
finite-enumeration paper-certificate candidate. A future certificate must
separate source, pattern, classifier, stability, and ledger fields. The
source field records the finite RNA alphabet, the two-axis readback, the
codon cube, and the standard genetic-code table. The pattern field records
local $Q_{2}$ faces, $3+1$ translation-class splits, corner partitions, and
alternative-code deformations. The classifier field records codon labels.
The stability field records closure of the finite enumeration under the
chosen readback. The ledger policy preserves the codon, six-bit display,
local face, free coordinates, translation label, and table provenance.
\end{definition}

\begin{definition}[RNA alphabet and window-six readback]
\label{def:philosophy-local-tile-rna-window-six-readback}
Let
$$
  \Sigma:=\{\mathrm{U},\mathrm{C},\mathrm{A},\mathrm{G}\}.
$$
The coordinate map $b:\Sigma\to\{0,1\}^{2}$ is
$$
\begin{aligned}
  b(\mathrm{U})&=00,&
  b(\mathrm{C})&=01,\\
  b(\mathrm{A})&=10,&
  b(\mathrm{G})&=11.
\end{aligned}
$$
For a codon $c=c_{1}c_{2}c_{3}\in\Sigma^{3}$, the six-bit readback is
$$
  \beta(c):=b(c_{1})\,b(c_{2})\,b(c_{3})\in\{0,1\}^{6}.
$$
Write $Q_{6}:=\{0,1\}^{6}$.
\end{definition}

\begin{definition}[Standard translation label map]
\label{def:philosophy-standard-translation-label-map}
The standard translation label map is the external biological table
$$
  \tau_{\mathsf{std}}:\Sigma^{3}\rightsquigarrow
  \mathsf{AminoAcidClass}\oplus\mathsf{Stop}.
$$
The rows used by this subsection are
$$
\begin{array}{c|l}
  \text{label} & \text{codons}\\
\hline
  \mathrm{Stop} & \mathrm{UAA},\mathrm{UAG},\mathrm{UGA}\\
  \mathrm{Trp}  & \mathrm{UGG}\\
  \mathrm{Ile}  & \mathrm{AUU},\mathrm{AUC},\mathrm{AUA}\\
  \mathrm{Met}  & \mathrm{AUG}\\
  \mathrm{Phe}  & \mathrm{UUU},\mathrm{UUC}\\
  \mathrm{Leu}  &
    \mathrm{UUA},\mathrm{UUG},\mathrm{CUU},\mathrm{CUC},
    \mathrm{CUA},\mathrm{CUG}.
\end{array}
$$
The biological labels are bridge data, not BEDC primitives.
\end{definition}

\begin{definition}[Axis-aligned local $Q_{2}$ face]
\label{def:philosophy-axis-aligned-local-q-two-face}
An axis-aligned local $Q_{2}$ face of $Q_{6}$ is obtained by choosing two
of the six bit positions to vary and fixing the other four. Such a face
has four vertices and is displayed as
$$
  F\simeq Q_{2}:=\{00,01,10,11\}.
$$
\end{definition}

\begin{definition}[Three-plus-one face]
\label{def:philosophy-three-plus-one-face}
An axis-aligned local $Q_{2}$ face $F$ is a $3+1$ face under a translation
label map $\tau$ when the multiset of labels on its four codons has
exactly two labels with multiplicities $3$ and $1$. The size-three label
class is the three-point class, and the remaining codon is the singleton
corner.
\end{definition}

\begin{definition}[Local two-bit tile and corner mechanisms]
\label{def:philosophy-local-two-bit-tile}
The local no-$11$ triangle is
$$
  X_{2}:=\{00,01,10\}
$$
inside
$$
  Q_{2}:=\{00,01,10,11\}.
$$
Thus
$$
  Q_{2}=X_{2}\sqcup\{11\}.
$$
An $11$-control tile is a $3+1$ face whose three-point class is $X_{2}$
and whose singleton corner is $11$. A $00$-singleton tile is a $3+1$ face
whose singleton corner is $00$ and whose three-point class is
$Q_{2}\setminus\{00\}$.
\end{definition}

\concretizedIn{ch:concrete-instances-dna-window-six-boundary-tile-namecert}{The local two-bit boundary-tile classifier is realized as finite codon-window, boundary-tile, transport, provenance, and NameCert rows.}

\begin{definition}[Distinguished local faces]
\label{def:philosophy-local-codon-q-two-face}
The two $11$-control faces are
$$
\begin{aligned}
  \mathsf{AUFace}
    &:=\{\mathrm{AUU},\mathrm{AUC},\mathrm{AUA},\mathrm{AUG}\},\\
  \mathsf{StopTrpFace}
    &:=\{\mathrm{UAA},\mathrm{UAG},\mathrm{UGA},\mathrm{UGG}\}.
\end{aligned}
$$
The two $00$-singleton faces are
$$
\begin{aligned}
  \mathsf{PheLeuFace}_{\mathrm{U}}
    &:=
      \{\mathrm{UUU},\mathrm{UUA},\mathrm{CUU},\mathrm{CUA}\},\\
  \mathsf{PheLeuFace}_{\mathrm{C}}
    &:=
      \{\mathrm{UUC},\mathrm{UUG},\mathrm{CUC},\mathrm{CUG}\}.
\end{aligned}
$$
\end{definition}

\begin{definition}[Reserved local corner]
\label{def:philosophy-reserved-local-corner}
For a local codon $Q_{2}$ face $F$, a reserved local corner is a singleton
corner whose biological bridge row is separated from the three-point
class. Reserved means singleton or control-sensitive under the translation
classifier. It does not mean absent, deleted, or identified with the BEDC
channel terminator.
\end{definition}

\begin{definition}[Stop graph and release-factor arms]
\label{def:philosophy-stop-graph-and-rf-arms}
On
$$
  X_{2}=\{00,01,10\}
$$
put the Hamming-distance-one graph
$$
  01\;-\;00\;-\;10.
$$
Under the stop-face readback this is
$$
  \mathrm{UAG}\;-\;\mathrm{UAA}\;-\;\mathrm{UGA}.
$$
The bacterial release-factor arms are external biological bridge rows:
$$
\begin{aligned}
  \mathsf{RF1Arm}&:=\{\mathrm{UAA},\mathrm{UAG}\},\\
  \mathsf{RF2Arm}&:=\{\mathrm{UAA},\mathrm{UGA}\}.
\end{aligned}
$$
\end{definition}

\begin{definition}[Local reassignment modes]
\label{def:philosophy-local-reassignment-modes}
A leaf reassignment on an $X_{2}\sqcup\{11\}$ tile moves one leaf of the
$X_{2}$ part to another translation class. The central displayed case is
$$
  \{00,01,10\}/\{11\}
  \rightsquigarrow
  \{00,01\}/\{10,11\}.
$$
An arm reassignment on the stop graph moves an edge such as
$\{00,01\}$ as one biological classifier row.
\end{definition}

\begin{definition}[Local translation state]
\label{def:philosophy-boundary-tile-finite-state-skeleton}
A local translation state has the schematic form
$$
  S_{t}:=(f_{t},w_{t},r_{t},\kappa_{t}),
$$
where $f_{t}$ is the frame, $w_{t}\in Q_{6}$ is the current window,
$r_{t}$ records release or ribosome state, and $\kappa_{t}$ records
cellular context. Under the standard-code bridge:
$$
\begin{gathered}
  w_{t}\notin\mathrm{Stop}
    \quad\Longrightarrow\quad S_{t}\to S_{t+1},\\
  w_{t}\in\mathrm{Stop}
    \quad\Longrightarrow\quad S_{t}\to S_{\mathsf{halt}}.
\end{gathered}
$$
This is a modeling interface, not a closed autonomous biological dynamics.
\end{definition}

\begin{lemma}[Codon cube and window-six bijection]
\label{lem:philosophy-codon-cube-window-six-bijection}
The readback $\beta$ gives a finite bijection
$$
  \Sigma^{3}\cong Q_{6}.
$$
\end{lemma}
\begin{proof}
Each base in $\Sigma$ has a unique two-bit coordinate under $b$, so a
three-base codon has a unique six-bit display. Conversely, every six-bit
word splits uniquely into three two-bit blocks, and each block is one of
$00$, $01$, $10$, and $11$, hence decodes uniquely to one of
$\mathrm{U}$, $\mathrm{C}$, $\mathrm{A}$, and $\mathrm{G}$.
\end{proof}

\begin{lemma}[Local two-bit partition]
\label{thm:philosophy-local-two-bit-partition}
The two-bit square satisfies
$$
  Q_{2}=X_{2}\sqcup\{11\},
  \qquad
  X_{2}=\{00,01,10\}.
$$
\end{lemma}
\begin{proof}
The four two-bit words are exactly $00$, $01$, $10$, and $11$. Removing
the word $11$ leaves exactly $00$, $01$, and $10$. The displayed union is
disjoint and exhaustive.
\end{proof}
